% Setup - do not change
\documentclass[11pt]{article}
\usepackage[top=0.9in, left=0.9in, bottom=0.9in, right=0.9in]{geometry} 
\usepackage{parskip}

\usepackage[english]{babel}
\usepackage[utf8]{inputenc}
\usepackage{amsmath,amsthm,amssymb,graphicx,pdfpages,lipsum,hyperref}
\usepackage[none]{hyphenat}
\usepackage{csquotes}

\setlength\parindent{0pt}
%%%%%%%%%%%%%%%%%%%%%%%%%%%%%%%%%%%%%%%%%%%%%%%%%%%%%%%%%%%%%%%%%%%
% add other packages here if required

%% Bibliography are specified in this file. You can also choose inline bib style if you want to. But make sure your citation style is consistent (and proper)
% For more details on citation: https://library.unimelb.edu.au/recite
\usepackage[sorting = none]{biblatex}
\addbibresource{references.bib}

%%%%%%%%%%%%%%%%%%%%%%%%%%%%%%%%%%%%%%%%%%%%%%%%%%%%%%%%%%%%%%%%%%% the '%' symbol denotes comments

% Begin document creation
% DELETE THE \lipsum PLACEHOLDERS WHEN YOU BEGIN
\title{Flight Arrivals as a Demand Signal:\\
Forecasting Yellow Taxi Revenue at JFK and LaGuardia}
\author{Tavish Balyan \\ Student ID: 1615873 \\
\href{https://github.com/MAST30034-AppliedDataScience/project-1-individual-SavvyHack/tree/951b6239bc4937ce3533f1b2ca4b42d81444ed06}{Github repo with commit}}
\date{August 31, 2026}
 
\begin{document}
\maketitle
 
\section{Introduction}
 
A yellow taxi driver joining the holding lot at JFK is betting that the queue
will clear before the fare stops being worth the wait, and the bet is a guess, but it doesn't need to be.
The arrivals that will fill the rank in two hours are already published on a
schedule that drivers do not consult. Can flight arrival data, combined with
weather, predict airport taxi demand accurately enough to make that decision a
calculated one? The audience is the individual yellow taxi driver choosing
between airport and city work at the start of a shift, who absorbs both a long
return and an unknown queue if they commit at the wrong hour; a useful prediction
must therefore estimate not only how many trips there will be but what each is
worth, per airport-hour.
 
We restrict the analysis to yellow taxis, as they alone may street-hail
passengers across all five boroughs, dominate the medallion ranks at both
airports \cite{tlc_userguide}, and operate under the flat JFK to Manhattan fare,
which makes their per-trip revenue structurally different from a for-hire
vehicle's and worth modelling separately. JFK (taxi zone 132) and LaGuardia
(zone 138) contrast usefully, as JFK is international, long-haul and flat-fared
whereas LGA is domestic, shorter-haul and metered; Newark enters only as a
context variable, as it sits outside the yellow medallion's service area.
 
The timeline is January 2023 to June 2024: 18 months of TLC Yellow Taxi records,
58,642,319 raw trips, of which 56,637,949 survive cleaning and 4,653,181 are
airport pickups. The window begins after post-pandemic ridership had stabilised,
as a model fitted across the recovery would learn a trend that no longer exists,
and 18 months spans two winters and two summers, the minimum needed to separate
weather from seasonal effects while leaving enough history for a forward split.
 
Two external datasets support the goal. Flight arrivals come from the Bureau of
Transportation Statistics Reporting Carrier On-Time Performance database
\cite{bts}, which records both scheduled and actual arrival times and so allows
two distinct signals to be built; one known to a driver in advance from the
published schedule, and one describing realised slippage observable only after
the fact. Hourly weather comes from Meteostat \cite{meteostat} for the JFK, LGA,
EWR and Central Park stations, on the reasoning that precipitation suppresses the
walk-up alternatives to a taxi and should therefore raise demand at the rank.
 
\textbf{Methodology overview.} The contribution is a two-model decomposition of
expected revenue per airport-hour. Trip records, flight arrivals and weather
observations are aggregated onto a common airport-hour spine of 26,256 rows.
Model 1 is a Poisson generalised linear model for the count of pickups in an
airport-hour and Model 2 a histogram-based gradient boosting regressor for the
mean fare of a trip in the same hour; their product is the quantity the queue
decision turns on. The two contrast deliberately, one parametric and
interpretable, the other free to capture interactions the first cannot express.
Section \ref{sec:results} verifies that their residuals are close to independent,
which is what licenses multiplying them.
 
\textbf{Assumptions.}
\begin{itemize}
  \item Every recorded pickup in zones 132 and 138 is assumed to serve an
        arriving passenger, as queue length is not observable in the data.
  \item \texttt{FlightDate} in the BTS extract is the departure date rather than
        the arrival date, so arrivals crossing midnight require a date shift.
  \item Cancelled flights are assumed to generate no taxi demand, and fare value
        is taken as the driver-relevant total excluding tips, as tipping is not
        predictable from an arrival schedule.
\end{itemize}
 
\section{Preprocessing}
 
\textbf{TLC Yellow Taxi records.} Eighteen monthly parquet files are loaded and
normalised to a common 19-column schema, as column naming drifts between
vintages. Records are then filtered against the business rules in the TLC Data
Dictionary \cite{tlc_dict}, with the surviving row count recorded at each step in
(Table \ref{tab:taxi_waterfall}).
 
\begin{table}[htb]
\centering
\footnotesize
\begin{tabular}{lrr}
\toprule
Filter & Rows remaining & Removed \\
\midrule
Raw ingest                              & 58,642,319 & --- \\
Pickup within study window              & 58,642,192 & 127 \\
Drop-off after pickup                   & 58,620,446 & 21,746 \\
Duration between 1 min and 6 h          & 57,922,347 & 698,099 \\
Positive trip distance                  & 57,243,639 & 678,708 \\
Positive fare and total                 & 56,642,200 & 601,439 \\
Metered fare at or above initial charge & 56,641,660 & 540 \\
Implied speed below 80 mph              & 56,637,949 & 3,711 \\
Documented rate code and payment type   & 56,637,949 & 0 \\
\bottomrule
\end{tabular}
\caption{Taxi cleaning waterfall; 2,004,370 records are discarded, or 3.42\% of
raw ingest. The final row removes nothing, which is the finding: both coded
fields are clean rather than unchecked.}
\label{tab:taxi_waterfall}
\end{table}
 
\begin{itemize}
  \item Restrict to pickups in zones 132 and 138: 4,653,181 airport pickups, or
        8.22\% of cleaned trips, split 2,753,043 JFK and 1,900,138 LGA.
  \item Aggregate to hourly counts and fare summaries: 25,076 non-empty
        airport-hours on a complete spine of 26,256; of the 1,180 remaining,
        1,175 are genuine zeros rather than missing data and five are
        daylight-saving artefacts.
\end{itemize}
 
\textbf{BTS Reporting Carrier On-Time Performance.} December 2022 to June 2024
loads as 1,349,639 rows, 14 columns after normalisation, and filters to JFK and
LGA arrivals: 461,205 rows (34.17\%), comprising 204,135 JFK and 257,070 LGA.
\begin{itemize}
  \item Convert the \texttt{hhmm} integer format, mapping the 2400 sentinel to 0
        before the overnight date shift that moves 4.79\% of arrivals to the
        following day, and bucket to the arrival hour only after both. The order
        matters, as reversing it misassigns every arrival crossing midnight.
  \item Drop two rows with implied block time above 18 hours (461,203 remaining),
        both with corrupt clock fields, then restrict to scheduled arrivals in
        the study window: 436,754, with 24,449 removed. The December 2022
        download is included because flights departing 31 December 2022 land in
        the early hours of 1 January 2023; it recovers 24 such arrivals, 4.7\% of
        that day's traffic.
  \item Of these, 10,389 were cancelled (2.38\%) and 1,634 diverted; with 61 more
        landing past the window close, 424,670 realised arrivals fall inside it,
        or 97.23\% of scheduled. Cancellations are excluded from realised counts,
        as retaining them would inflate predicted demand on exactly the disrupted
        days where prediction matters most. Aggregation gives 26,256
        airport-hours, plus 13,128 Newark hours as context.
\end{itemize}
 
\textbf{Meteostat hourly observations.} Four stations return 52,047 rows against
a complete grid of 52,512 site-hours, near-complete at the airports (two missing
hours each) but sparse at Central Park (459 missing).
\begin{itemize}
  \item Drop the \texttt{coco} condition code, 93--97\% null by station and
        fog-only among survivors. Impute \texttt{prcp} from a neighbouring
        station, or dry where none reported, affecting 1,527 airport-hours
        (5.8\%); these and four hours carried forward across the spring-forward
        gap are the only imputations in the pipeline.
  \item Derive \texttt{is\_freezing} from $\texttt{prcp} > 0$ and
        $\texttt{temp} \le 0^\circ$C, flagging 0.32\% of hours. Snow is not used,
        as the window coincides with a record snow drought.
\end{itemize}
 
\textbf{Join.} All sources are left-joined onto the 26,256-row spine with row
counts asserted after every join, so any silent fan-out is caught where it
occurs. The table grows from 20 columns to 31 with JFK and LGA flights, 36 with
Newark context, 47 with weather and 49 with holiday flags.
 
\section{Analysis and Geospatial Visualisation}
\label{sec:analysis}
 
Every distribution below is computed over the full cleaned table, as percentiles
and outlier fences are exactly the quantities a sample would distort; only the
scatter in (Figure \ref{fig:fare}b) is subsampled, at 1\%, as 4.65M points render
as a solid block. No trip-level value is imputed anywhere: the weather fills of
Section 2 touch 5.8\% of airport-hours and none of the response.
 
\begin{figure}[htb]
\centering
\includegraphics[width=0.84\textwidth]{fig01_fare_distribution.png}
\caption{Metered fare distribution, citywide against airport pickups. The JFK
mass point at \$70.00 is the published flat fare to Manhattan.}
\label{fig:fare}
\end{figure}
 
From (Figure \ref{fig:fare}) the airport fare distribution is not a shifted copy
of the citywide one but a different shape: JFK is bimodal, with 1,492,107 pickups
(54.2\% of JFK trips) billing exactly \$70.00, whereas LGA is unimodal and
metered. A mean fare at JFK is therefore a mixing weight between two regimes
rather than a central tendency, which is why Model 2 is given a learner that can
represent a mixture.
Applying Tukey's fence \cite{tukey} to the full distribution, 66,018 airport
trips (1.42\%) exceed the \$114.85 upper bound; profiling by rate code shows
50.95\% are the Nassau or Westchester out-of-city rate and 22.26\% negotiated
fares, both documented regimes \cite{tlc_dict}, so these records are extreme
rather than invalid and are retained. One record carries 67.1\% of the variance
of citywide \texttt{fare\_amount}, so quartiles rather than standard deviations
are quoted throughout.
 
The target needs the same scrutiny before a family is chosen. Of the 26,256
airport-hours, 1,180 record no pickups at all, split 1,133 LGA against 47 JFK, so
LGA closes overnight while JFK merely thins. The variance-to-mean ratio computed
\emph{within} (airport, hour-of-day) cells exceeds one in all 48 cells, at medians
of 12.44 for JFK and 21.40 for LGA; the conditional test is the relevant one, as
an unconditional figure would mostly measure the gap between 03:00 and 18:00.
That overdispersion survives conditioning is what makes the negative binomial
comparison in Section \ref{sec:results} a genuine question rather than a
formality. Volume and value are also coupled unequally: across the 168 (day,
hour) cells the Spearman correlation between mean pickups and mean fare is
$+0.771$ at JFK but only $+0.536$ at LGA, so LaGuardia has hours that pay near
the top of its range on a short queue --- a gap Section \ref{sec:recs} turns into
a shift plan.
 
\begin{figure}[htb]
\centering
\includegraphics[width=0.84\textwidth]{fig04_flights_to_pickups.png}
\caption{Correlation between hourly pickups and scheduled arrivals at lags of
$-2$ to $+2$ hours, and domestic-arrival coverage by hour of day.}
\label{fig:flights}
\end{figure}
 
From (Figure \ref{fig:flights}) the correlation between pickups and scheduled
arrivals peaks at a one-hour lag at both airports, 0.717 at JFK and 0.842 at LGA,
falling away either side; this justifies carrying
\texttt{sched\_arrivals\_prev\_hr} into the feature set, as an arriving passenger
clears the terminal before reaching the rank. The gap between the airports is a
coverage artefact rather than a behavioural one: JFK converts 14.25 pickups per
domestic arrival against LGA's 7.80, but 38.5\% of JFK's scheduled domestic
arrivals are long-haul against 3.4\% at LGA, and BTS reports domestic carriers
only. The flight feed therefore sees most of LaGuardia and roughly two thirds of
JFK, which caps how well any schedule signal can do at JFK.
 
Weather is the clearest negative result. Matched within (airport, hour-of-day)
cells, wet hours carry 7.0\% fewer pickups at LGA (paired $t = -2.39$,
$p = 0.025$) and no significant change at JFK, so rain does not raise airport
demand as the introduction supposed; the walk-up alternatives that rain
suppresses in Manhattan barely exist at a rank where the queue is already the
default.
 
% Cropped in LaTeX, not re-plotted. The 2x2 grid has an all-white gutter whose
% midpoint sits at exactly 50% of the image height, so trimming the bottom half
% keeps panels (a) and (b) and discards the mean-fare map and the unreferenced
% "Night minus day" panel. The trim is given in absolute big points because
% \height is not defined inside graphicx's trim key: the image is 1425px at
% 300dpi = 342bp tall, so half is 171bp.
% Width 0.62: the panels carry large filled polygons rather than fine tick
% labels, so this is the figure that gives up the most page for the least.
\begin{figure}[htb]
\centering
\includegraphics[trim={0 171bp 0 0}, clip,
                 width=0.62\textwidth]{fig06_dropoff_geography.png}
\caption{Share of each airport's drop-offs by taxi zone, on a log scale. A
further 3.16\% of airport trips end in zones 264/265, which have no geometry and
are not mapped.}
\label{fig:geo}
\end{figure}
 
Pickups occur in exactly two zones, so the informative map is of where airport
trips end. From (Figure \ref{fig:geo}) the 4,653,181 airport pickups disperse
across 261 drop-off zones, 60.81\% terminating in Manhattan at a mean of \$82.48
against 18.00\% in Queens at \$46.72. The two panels are near-indistinguishable,
which is itself the finding: destination mix is a property of the city rather
than of the airport, so a driver gains nothing by choosing a rank on where its
passengers are headed.
 
What the map conceals is that fare is only half of what an hour is worth. Times
Sq/Theatre District is the largest single destination at 5.38\% and pays \$86.28,
but the run takes 49.7 minutes and returns only \$104.16 per engaged hour,
against Murray Hill's \$136.12 on \$82.03 over 36.2 minutes. Street-hail rate is
a second and separate axis: Times Square supplies 144.3 citywide pickups an hour
against 69.7 in the Garment District, so how fast this trip ends and where the
next one begins are different questions, and Section \ref{sec:recs} needs both.
 
\section{Modelling}
\label{sec:modelling}
 
Expected revenue in an airport-hour is decomposed into two questions a driver
asks separately, as how many fares are waiting and what each is worth are
governed by different mechanisms:
\begin{equation}
\widehat{R}_{a,h} = \widehat{N}_{a,h} \times \widehat{F}_{a,h},
\label{eq:product}
\end{equation}
where $\widehat{N}_{a,h}$ is a Poisson prediction of pickup count and
$\widehat{F}_{a,h}$ a gradient-boosted prediction of mean fare per trip.
 
Both models train on 2023 and test on January to June 2024, a strictly forward
split, as any recommendation to a driver must be built only from information
available before the shift begins. Of the 26,256 airport-hours, 342 are excluded
upstream (six daylight-saving anomalies, 336 lacking a full weekly lag), leaving
17,180 training and 8,734 test rows; Model 2 is defined only where at least one
trip occurred, so it sees 16,375 and 8,375. A roles manifest bars 23 columns so
that no same-hour outcome reaches either feature list, standardisation is fitted
on the training split alone, and \texttt{arrival\_slippage\_lag\_1h} is dropped as
an exact linear combination of two retained features. Everything below is
reported beside a seasonal-naive benchmark \cite{hyndman} predicting each hour
with the same hour one week earlier, as a model that cannot beat last Tuesday
adds nothing a driver does not already have. Hour enters as 23 dummies
interacted with airport rather than as a continuous term, as demand is not
monotone.
 
\textbf{Model 1: a Poisson GLM for demand.} The response is a non-negative
integer count with a hard floor at zero, so a Poisson GLM with a log link is the
natural specification \cite{mccullagh}, and its multiplicative coefficients read
directly as the proportional effect of an extra scheduled arrival. The design
matrix is of full rank 98, so every coefficient is identified; the Pearson
dispersion, however, is 10.48 rather than the assumed 1, consistent with Section
\ref{sec:analysis}, so a quasi-Poisson scale \cite{wedderburn} is applied,
inflating the standard errors 3.24 times, and no coefficient is interpreted
without it. Overdispersion of that size invites a negative binomial alternative,
which assumes a different variance shape rather than merely rescaling the errors;
it is fitted on the same design, so the two count models differ in exactly one
assumption.
 
\textbf{Model 2: gradient-boosted trees for the fare.} A contrasting family is
used deliberately rather than a second GLM. The fare at an airport-hour is a
mixture of flat-rate and metered trips, as Section \ref{sec:analysis}
established, so its conditional mean is a non-linear function that a tree
ensemble can capture without the interactions being specified in advance
\cite{friedman, sklearn}. Hyperparameters are tuned on a forward validation
split, training to 30 September 2023 and validating on the final quarter, so the
configuration is chosen by a candidate with no holiday history; the shortfall is
identical across the grid, so the ranking holds and the refit uses the whole year.
 
\section{Results and Discussion}
\label{sec:results}
 
\begin{table}[htb]
\centering
\footnotesize
\begin{tabular}{llrrr}
\toprule
Target & Model & $n$ & RMSE & MAE \\
\midrule
Pickups per hour    & Seasonal naive (benchmark)  & 8,734 & 52.84 & 36.08 \\
                    & Poisson GLM                 & 8,734 & \textbf{40.64} & \textbf{28.89} \\
                    & Negative binomial GLM       & 8,734 & 44.09 & 30.93 \\
\addlinespace
Mean fare (US\$)    & Seasonal naive (benchmark)  & 8,198 & 10.13 & 5.18 \\
                    & Boosted trees               & 8,198 & \textbf{7.11} & \textbf{3.77} \\
                    & Boosted trees, full support & 8,375 & 8.43 & 4.13 \\
\addlinespace
Revenue (US\$/hour) & Benchmark product           & 8,734 & 4,175 & 2,834 \\
                    & Model 1 $\times$ Model 2    & 8,734 & \textbf{3,230} & \textbf{2,255} \\
\bottomrule
\end{tabular}
\caption{Test-period losses, January to June 2024. Rows compare only within a
target and at equal $n$: the value benchmark is undefined in the 177 hours whose
counterpart last week was empty. The count models are also compared on mean
Poisson deviance, 10.021 against 10.885 in the Poisson's favour, since squared
error alone favours the Poisson by construction.}
\label{tab:results}
\end{table}
 
In sample the negative binomial wins decisively on AIC, 182,996 against 280,978,
but the verdict reverses out of sample, where it is worse on all three losses in
(Table \ref{tab:results}). The in-sample criterion rewards a dispersion parameter
fitted to 2023, whereas the test period asks whether that flexibility generalises;
it does not, the Poisson is retained, and both beat the benchmark.
 
For the fare model the improvement on common support is 29.8\% in RMSE. The
benchmark cannot score 177 hours at all, as their counterpart last week was
empty, and Model 2 scores those at \$32.02; excluding them would flatter the
comparison, so the honest figure is the \$8.43 over full support.
 
% Not cropped: panel (a)'s plot area overlaps where panel (b)'s tick labels
% begin, so no clean vertical cut exists.
\begin{figure}[htb]
\centering
\includegraphics[width=0.80\textwidth]{fig09_model_interpretation.png}
\caption{Model 1 fitted demand across a midweek June day, and Model 2
permutation importance in dollars of test RMSE added when a feature is shuffled.}
\label{fig:interp}
\end{figure}
 
Permutation importance produces the most consequential negative result (Figure
\ref{fig:interp}). The fare model is dominated by its own history, with
\texttt{mean\_total\_same\_hour\_7d} contributing \$1.257 of RMSE against \$0.095
for the three collinear flight columns shuffled jointly. Flight schedules predict
how many passengers will want a taxi but almost nothing about what each trip will
be worth, as the schedule sets arrival volume while fare is set by destination;
the external dataset earns its place through Model 1 rather than Model 2.
 
Model 1's rate ratios say what that place is. A further standard deviation of the
current hour's scheduled arrivals multiplies expected pickups by 1.193 (95\% CI
1.164--1.224), and of realised arrivals in the previous hour by 1.150
(1.137--1.164). The lagged schedule term itself reads 0.918 (0.869--0.969),
against the positive correlation of Section \ref{sec:analysis}; this is
collinearity rather than contradiction, as without the slippage difference
\texttt{sched\_arrivals\_prev\_hr} and \texttt{actual\_arrivals\_lag\_1h} are near
duplicates and the fit splits one effect between them, the realised count
absorbing it. The informative coefficient is \texttt{share\_delayed\_15\_lag\_1h}
at 1.055 (1.048--1.061); holding arrival volume fixed, an hour in which more
flights ran late is followed by \emph{more} pickups, not fewer. Delay displaces
demand forward rather than destroying it, which is the opposite of what a driver
watching a departures board would infer.
 
Model 1 under-predicts by 6.8 pickups per hour, with 71\% of the 48 airport-hour
cells sharing that sign, and its structure matters more than its magnitude. It is
not a year-on-year level effect, as demand fell 9.8\% at JFK and 3.9\% at LGA
between matched months while the residuals point the other way, nor monotone in
the level: it peaks in the seventh decile at +20.4 pickups and falls to near zero
at both extremes, so the model is conservative in the moderately busy middle
where a shift is mostly spent. The bias also decays across the horizon, from
+13.7 pickups in the first test month to +0.2 in the sixth; this is a handover
effect rather than a standing property, as the weekly lags entering January still
point back at the holiday lull and progressively stop doing so, which makes the
caveat time-limited.
 
Multiplying two predictions is legitimate only if their errors are close to
independent, as any residual covariance is a term Equation \ref{eq:product}
omits. That term is \$24 per airport-hour against a mean observed revenue of
\$12,986, or 0.2\% of level, so the product is sound. Over the test period the
combination improves on the benchmark product by 22.6\% and totals \$107.5M
predicted against \$113.4M observed; that 5.2\% shortfall is the early-horizon
conservatism above rather than a defect of the combination. Refitting Model 2
with hours weighted by pickup count moved combined RMSE by 0.2\%, so the
unweighted fit is retained.
 
\section{Recommendations}
\label{sec:recs}
 
\begin{figure}[htb]
\centering
\includegraphics[width=0.84\textwidth]{fig10_expected_revenue.png}
\caption{Predicted expected revenue per airport-hour across an ordinary June
midweek day, JFK against LGA, with observed values overlaid, and the JFK
advantage across the week.}
\label{fig:revenue}
\end{figure}
 
\textbf{Read the hour behind you, and treat a bank of delays as a reason to
stay.} The instinct when the board fills with delays is that the rank has gone
dead, so a driver leaves for the city. The model says otherwise: with arrival
volume held fixed, a standard deviation more of
\texttt{share\_delayed\_15\_lag\_1h} multiplies the next hour's pickups by 1.055,
realised arrivals in the previous hour carry a rate ratio above one, and (Figure
\ref{fig:flights}) puts the peak correlation at a one-hour lag. Delays move
passengers later; they do not remove them. The operational form needs nothing but
a phone: read the previous hour's landings rather than the current hour's
schedule, and discount a thin next hour if the hour behind it ran late. The
qualification is seasonal, as Section \ref{sec:results}'s early-horizon bias is
concentrated in January, so the rule should be trusted least in the first weeks
of a year and the model refitted monthly.
 
\textbf{Start a weekday at LaGuardia and move to JFK after midday.} JFK is larger
and pays more, so the default is to work it all day, and for most of the clock
that is right: predicted revenue peaks at JFK 16:00 at \$28,214 per hour (the
amount the fleet earned collectively), and the five best JFK hours fall between
15:00 and 22:00 at fares near \$85. The exception is narrow but repeatable: from
(Figure \ref{fig:revenue}) LGA out-earns JFK in 19 of the 168 day-hour cells,
concentrated at 09:00 to 11:00 on weekdays and peaking on Wednesday at 10:00 with
an advantage of \$6,668 per hour. Section
\ref{sec:analysis} explains why this stays invisible: LaGuardia couples volume to
value at only $+0.536$ against JFK's $+0.771$, so those cells pay near the top of
the range on a short queue, and a driver judging by queue length reads them as
quiet. The advantage is strictly midweek. For a Queens-based driver the plan is a
weekday start at LGA through the morning bank and a move to JFK from midday; for
one already at JFK it is a reason not to rejoin the lot at 09:00.
 
\textbf{Judge the drop-off by the hour it returns and the street-hail rate it
lands in, not by the fare.} A driver has no control over the destination but
does control what follows it, and Section \ref{sec:analysis} shows that two
quantities decide it: how fast the trip ends, and how easily the next one
starts. Upper East Side North and
Midtown East return \$139.22 and \$127.60 per engaged hour into zones supplying
184.7 and 153.1 citywide pickups an hour, so after either the street is the right
follow-up and a return to the rank wastes the trip back. Times Square is the
opposite case on the higher fare, at \$104.16 per engaged hour, and the Garment
District worse again on street-hail rate at 69.7 an hour; after those, (Figure
\ref{fig:revenue}) decides whether the return leg pays at that hour --- which,
for a Garment District drop between 15:00 and 22:00, it does. Ranking the four or
five zones a driver actually works is a one-off exercise.
 
\section{Conclusion}
 
% \looseness asks TeX to set this paragraph one line shorter by tightening
% inter-word spacing within its normal tolerance. No content is lost and no
% length is redefined.
{\looseness=-1
Published flight schedules predict airport taxi demand, but only demand: Model 1
beats a seasonal-naive benchmark by 23.1\% and the combination by 22.6\% on
revenue, while arrival data says almost nothing about fare, which is set by
destination. The value lies where the schedule contradicts intuition: delays
defer demand rather than destroy it, and a midweek LaGuardia morning out-earns
JFK. Both are actionable from a phone at the rank.\par}
 
\printbibliography
 
\end{document}